\section{خواص سری فوریه، قضیۀ پارسوال و صورت مختلط}
\def\EMfootchapter{فصل ۲ \ \lr{\textbullet}\ خواص سری فوریه، قضیۀ پارسوال و صورت مختلط}
\def\EMtitlesize{26}\def\EMtitleskip{31}
\begin{frame}{فصل ۲: خواص سری فوریه، قضیۀ پارسوال و صورت مختلط}
{\fontsize{20}{26}\selectfont\bfseries\color{SLdef}\EMalign مطالب این فصل\par}
\vspace{10bp}
\EMitem{\lr{\textbullet}}{\dimexpr\linewidth-26bp\relax}{تقریب یک تابع متناوب توسط یک سری مثلثاتی با تعداد جملات محدود}
\EMitem{\lr{\textbullet}}{\dimexpr\linewidth-26bp\relax}{مشتق و انتگرال سری فوریه}
\EMitem{\lr{\textbullet}}{\dimexpr\linewidth-26bp\relax}{قضیۀ رایلی}
\EMitem{\lr{\textbullet}}{\dimexpr\linewidth-26bp\relax}{قضیۀ پارسوال}
\EMitem{\lr{\textbullet}}{\dimexpr\linewidth-26bp\relax}{صورت مختلط سری فوریه}
\end{frame}

\def\EMtitlesize{26}\def\EMtitleskip{31}
\begin{frame}{تقریب یک تابع متناوب توسط یک سری مثلثاتی با تعداد جملات محدود}
\begin{EMbox}{thm}{1}{}
\begin{uncoverenv}<2->
\EMp{برای تقریب \EMim{N} جمله‌ای یک تابع متناوب \EMim{f(x)} توسط یک سری سینوسی و کسینوسی، بهترین ضرایب (با معیار \lr{MMSE}) ضرایب سری فوریه می‌باشند.}
\end{uncoverenv}
\end{EMbox}
\begin{uncoverenv}<3->
\EMp{تابع متناوب با دورۀ تناوب \EMim{2\pi} را در نظر بگیرید:}
\vspace{3bp}
\EMeq{f(x) =a_0+\sum_{n=1}^{\infty}a_n\cos nx+b_n\sin nx,\qquad f(x) =f(x+2\pi)}
\end{uncoverenv}
\EMgap
\begin{uncoverenv}<4->
\EMp{تقریب \EMim{N} جمله‌ای مثلثاتی:}
\vspace{3bp}
\EMeq{f(x) \approx a_0+\sum_{n=1}^{N}a_n\cos nx+b_n\sin nx}
\end{uncoverenv}
\end{frame}

\def\EMtitlesize{26}\def\EMtitleskip{31}
\begin{frame}{تقریب یک تابع متناوب توسط یک سری مثلثاتی با تعداد جملات محدود (ادامه)}
\begin{uncoverenv}<2->
\EMp{آیا تقریب مثلثاتی مقابل \textbf{بهتر} نیست؟}
\vspace{3bp}
\EMeq{F(x)\triangleq\alpha_0+\sum_{n=1}^{N}\alpha_n\cos nx+\beta_n\sin nx}
\end{uncoverenv}
\EMgap
\begin{uncoverenv}<3->
\EMp{\textbf{معیار} بهتر بودن تقریب چیست؟ یک معیار: تابع خطا:}
\vspace{3bp}
\EMeq{E(x) =\big|f(x)-F(x)\big| =\left|f(x)-\alpha_0-\sum_{n=1}^{N}\big(\alpha_n\cos nx+\beta_n\sin nx\big)\right|}
\end{uncoverenv}
\EMgap
\begin{uncoverenv}<4->
\EMp{یک معیار خوب: متوسط مجذور خطا:}
\vspace{3bp}
\EMeq{D\triangleq\int_{-\pi}^{\pi}E^2(x)\,dx =\int_{-\pi}^{\pi}\left|f(x)-\alpha_0-\sum_{n=1}^{N}\big(\alpha_n\cos nx+\beta_n\sin nx\big)\right|^2dx}
\end{uncoverenv}
\end{frame}

\def\EMtitlesize{26}\def\EMtitleskip{31}
\begin{frame}{تقریب یک تابع متناوب توسط یک سری مثلثاتی با تعداد جملات محدود (ادامه)}
\begin{EMbox}{prf}{1}{}
\begin{uncoverenv}<2->
\EMp{برای بهترین تقریب متوسط مجذور خطا را کمینه می‌کنیم. این کار با مشتق‌گیری از \EMim{D} بر حسب ضرایب سری مثلثاتی و برابر قرار دادن این مشتق با صفر به دست می‌آید. به عبارت دیگر:}
\vspace{3bp}
\EMeq{\frac{\partial D}{\partial\alpha_0} =-2\int_{-\pi}^{\pi}\left[f(x)-\alpha_0-\sum_{n=1}^{N}\big(\alpha_n\cos nx+\beta_n\sin nx\big)\right]dx =0}
\end{uncoverenv}
\EMgap
\begin{uncoverenv}<3->
\EMeq{\int_{-\pi}^{\pi}f(x)\,dx =\int_{-\pi}^{\pi}\left[\alpha_0+\sum_{n=1}^{N}\big(\alpha_n\cos nx+\beta_n\sin nx\big)\right]dx =2\pi\alpha_0\quad \Longrightarrow\quad \alpha_0 =\frac1{2\pi}\int_{-\pi}^{\pi}f(x)\,dx =a_0}
\end{uncoverenv}
\EMgap
\begin{uncoverenv}<4->
\EMeq{\frac{\partial D}{\partial\alpha_m} =0\quad \Longrightarrow\quad \alpha_m =\frac1\pi\int_{-\pi}^{\pi}f(x)\cos mx\,dx =a_m}
\end{uncoverenv}
\EMgap
\begin{uncoverenv}<5->
\EMeq{\frac{\partial D}{\partial\beta_m} =0\quad \Longrightarrow\quad \beta_m =\frac1\pi\int_{-\pi}^{\pi}f(x)\sin mx\,dx =b_m}
\end{uncoverenv}
\EMgap
\begin{uncoverenv}<6->
\EMp{بهترین تقریب مثلثاتی با تعداد جملات محدود از یک تابع تناوبی سری فوریه با جملات محدود است.}
\end{uncoverenv}
\end{EMbox}
\end{frame}

\def\EMtitlesize{26}\def\EMtitleskip{31}
\begin{frame}{مشتق و انتگرال سری فوریه}
\begin{EMbox}{thm}{1}{}
\begin{uncoverenv}<2->
\EMp{اگر تابع \EMim{f(x)} با دورۀ تناوب \EMim{T} تناوبی بوده، پیوسته و دارای سری فوریه باشد، آنگاه مشتق این تابع نیز تناوبی بوده سری فوریه آن با استفاده از مشتق‌گیری از سری فوریه تابع به دست می‌آید.}
\vspace{3bp}
\EMeq{f(x) =a_0+\sum_{n=1}^{\infty}a_n\cos\frac{2\pi}Tnx+b_n\sin\frac{2\pi}Tnx}
\end{uncoverenv}
\EMgap
\begin{uncoverenv}<3->
\EMeq{\Longrightarrow\quad f'(x) =\sum_{n=1}^{\infty}\left(\frac{2\pi}Tnb_n\right)\cos\frac{2\pi}Tnx+\left(-\frac{2\pi}Tna_n\right)\sin\frac{2\pi}Tnx}
\end{uncoverenv}
\end{EMbox}
\begin{EMbox}{thm}{3}{}
\begin{uncoverenv}<4->
\EMp{اگر تابع \EMim{f(x)} با دورۀ تناوب \EMim{T} تناوبی بوده و دارای سری فوریه باشد، آنگاه انتگرال معین این تابع با استفاده از انتگرال‌گیری از سری فوریه تابع قابل محاسبه است.}
\vspace{3bp}
\EMeq{f(x) =a_0+\sum_{n=1}^{\infty}a_n\cos\frac{2\pi}Tnx+b_n\sin\frac{2\pi}Tnx}
\end{uncoverenv}
\end{EMbox}
\end{frame}

\def\EMtitlesize{26}\def\EMtitleskip{31}
\begin{frame}{مشتق و انتگرال سری فوریه (ادامه)}
\begin{EMbox}{thm}{1}{}
\begin{uncoverenv}<2->
\EMeq{\Longrightarrow\quad \int_0^xf(t)\,dt =a_0x+\sum_{n=1}^{\infty}\left(\frac T{2\pi n}a_n\sin\frac{2\pi}Tnx-\frac T{2\pi n}b_n\cos\frac{2\pi}Tnx+\frac T{2\pi n}b_n\right)}
\end{uncoverenv}
\end{EMbox}
\end{frame}

\def\EMtitlesize{26}\def\EMtitleskip{31}
\begin{frame}{قضیۀ رایلی}
\begin{EMbox}{thm}{1}{}
\begin{uncoverenv}<2->
\EMp{اگر دو تابع \EMim{f(x)} و \EMim{g(x)} با دورۀ تناوب \EMim{T} تناوبی بوده و سری فوریه این دو تابع به صورت زیر وجود داشته باشند:}
\vspace{3bp}
\EMeq{f(x) =a_0+\sum_{n=1}^{\infty}a_n\cos\frac{2\pi}Tnx+b_n\sin\frac{2\pi}Tnx}
\end{uncoverenv}
\EMgap
\begin{uncoverenv}<3->
\EMeq{g(x) =\alpha_0+\sum_{n=1}^{\infty}\alpha_n\cos\frac{2\pi}Tnx+\beta_n\sin\frac{2\pi}Tnx}
\end{uncoverenv}
\EMgap
\begin{uncoverenv}<4->
\EMp{آنگاه خواهیم داشت:}
\vspace{3bp}
\EMeq{\frac2T\int_Tf(x)\,\overline{g(x)}\,dx =2a_0\overline{\alpha_0}+\sum_{n=1}^{\infty}a_n\overline{\alpha_n}+b_n\overline{\beta_n}}
\end{uncoverenv}
\end{EMbox}
\end{frame}

\def\EMtitlesize{26}\def\EMtitleskip{31}
\begin{frame}{قضیۀ رایلی (ادامه)}
\begin{EMbox}{prf}{1}{}
\begin{uncoverenv}<2->
\EMp{اثبات با فرض حقیقی بودن \EMim{g(x)}:}
\vspace{3bp}
\EMeq{\begin{aligned}
\frac2T\int_Tf(x)g(x)\,dx&=\frac2T\int_Tf(x)\left(\alpha_0+\sum_{n=1}^{\infty}\alpha_n\cos\frac{2\pi}Tnx+\beta_n\sin\frac{2\pi}Tnx\right)dx\\
&=\frac2T\int_Tf(x)\alpha_0\,dx+\sum_{n=1}^{\infty}\frac2T\int_Tf(x)\left(\alpha_n\cos\frac{2\pi}Tnx+\beta_n\sin\frac{2\pi}Tnx\right)dx\\
&=\underbrace{\frac2T\int_Tf(x)\,dx}_{2a_0}\,\alpha_0+\sum_{n=1}^{\infty}\underbrace{\frac2T\int_Tf(x)\cos\frac{2\pi}Tnx\,dx}_{a_n}\,\alpha_n+\underbrace{\frac2T\int_Tf(x)\sin\frac{2\pi}Tnx\,dx}_{b_n}\,\beta_n\\
&=2a_0\alpha_0+\sum_{n=1}^{\infty}a_n\alpha_n+b_n\beta_n
\end{aligned}}
\end{uncoverenv}
\end{EMbox}
\end{frame}

\def\EMtitlesize{26}\def\EMtitleskip{31}
\begin{frame}{قضیۀ پارسوال}
\begin{EMbox}{thm}{1}{}
\begin{uncoverenv}<2->
\EMp{اگر تابع \EMim{f(x)} با دورۀ تناوب \EMim{T} تناوبی بوده و سری فوریه آن به صورت زیر وجود داشته باشد:}
\vspace{3bp}
\EMeq{f(x) =a_0+\sum_{n=1}^{\infty}a_n\cos\frac{2\pi}Tnx+b_n\sin\frac{2\pi}Tnx}
\end{uncoverenv}
\EMgap
\begin{uncoverenv}<3->
\EMp{آنگاه خواهیم داشت:}
\vspace{3bp}
\EMeq{\frac2T\int_T|f(x)|^2\,dx =2|a_0|^2+\sum_{n=1}^{\infty}\big(|a_n|^2+|b_n|^2\big)}
\end{uncoverenv}
\end{EMbox}
\begin{EMbox}{prf}{3}{}
\begin{uncoverenv}<4->
\EMp{در تساوی رایلی اگر به جای \EMim{g(x)} تابع \EMim{f(x)} را قرار دهیم به تساوی پارسوال می‌رسیم.}
\end{uncoverenv}
\end{EMbox}
\begin{uncoverenv}<5->
\EMp{با توجه به تساوی پارسوال مقدار توان تابع \EMim{f(x)} را می‌شود با استفاده از ضرایب سری فوریه به دست آورد. اگر تابع توسط تقریب \EMim{N} جمله‌ای از سینوس‌ها و کسینوس‌ها بیان شود طبق رابطۀ بالا توان تابع تقریبی کمتر از توان تابع \EMim{f(x)} خواهد بود و داریم:}
\end{uncoverenv}
\end{frame}

\def\EMtitlesize{26}\def\EMtitleskip{31}
\begin{frame}{قضیۀ پارسوال (ادامه)}
\begin{EMbox}{imp}{1}{}
\begin{uncoverenv}<2->
\EMeq{\frac2T\int_T|f(x)|^2\,dx\ \ge\ 2|a_0|^2+\sum_{n=1}^{N}\big(|a_n|^2+|b_n|^2\big)}
\end{uncoverenv}
\end{EMbox}
\end{frame}

\def\EMtitlesize{26}\def\EMtitleskip{31}
\begin{frame}{قضیۀ پارسوال (ادامه)}
\begin{EMbox}{ex}{1}{}
\EMp{برای تابع \EMim{f(x)} زیر، چند جمله از سری فوریه نیاز است تا حداقل ۹۴ درصد توان در تابع تقریبی باقی بماند.}
\EMgap
\begin{uncoverenv}<2->
\EMeq{f(x) =\sum_{n=1}^{\infty}\frac{2k}{n\pi}(1-\cos n\pi)\sin nx}
\end{uncoverenv}
\EMgap
\begin{uncoverenv}<3->
\EMeq{\frac2T\int_T|f(x)|^2\,dx =\frac1\pi\int_{-\pi}^{\pi}k^2\,dx =\frac2\pi\big[k^2x\big]_0^\pi =2k^2}
\end{uncoverenv}
\EMgap
\begin{uncoverenv}<4->
\EMeq{\frac{\displaystyle\sum_{n=1}^{N}\left(\frac{2k}{n\pi}(1-\cos n\pi)\right)^2}{2k^2}\ge0.94}
\end{uncoverenv}
\end{EMbox}
\end{frame}

\def\EMtitlesize{26}\def\EMtitleskip{31}
\begin{frame}{قضیۀ پارسوال (ادامه)}
\begin{minipage}[t]{500bp}\vspace{0pt}
\begin{EMbox}{ex}{1}{}
\begin{uncoverenv}<2->
\EMeq{\frac{\displaystyle\sum_{n=1}^{7}\left(\frac{2k}{n\pi}(1-\cos n\pi)\right)^2}{2k^2} \cong0.9496,\qquad \\[4pt] \frac{\displaystyle\sum_{n=1}^{5}\left(\frac{2k}{n\pi}(1-\cos n\pi)\right)^2}{2k^2} \cong0.9331\quad \Longrightarrow\quad N =7}
\end{uncoverenv}
\end{EMbox}
\end{minipage}
\EMside{237.1bp}{\uncover<3->{\EMfigure{1.121}{m02-power-example}{موج مربعی \EMim{\pm k} با دورۀ تناوب \EMim{2\pi} در مثال توان}}}
\end{frame}

\def\EMtitlesize{26}\def\EMtitleskip{31}
\begin{frame}{صورت مختلط سری فوریه}
\begin{uncoverenv}<2->
\EMp{فرض کنید تابع \EMim{f(x)} تناوبی و صرفاً جهت سادگی دارای دورۀ تناوب \EMim{2\pi} می‌باشد. لذا داریم:}
\vspace{3bp}
\EMeq{f(x) =a_0+\sum_{n=1}^{\infty}a_n\cos nx+b_n\sin nx}
\end{uncoverenv}
\EMgap
\begin{uncoverenv}<3->
\EMp{با استفاده از رابطۀ اویلر \EMim{e^{ix}=\cos x+i\sin x}:}
\vspace{3bp}
\EMeq{\cos nx =\frac{e^{inx}+e^{-inx}}2,\qquad \sin nx =\frac{e^{inx}-e^{-inx}}{2i},\qquad \frac1i =-i}
\end{uncoverenv}
\EMgap
\begin{uncoverenv}<4->
\EMeq{f(x) =a_0+\sum_{n=1}^{\infty}a_n\left(\frac{e^{inx}+e^{-inx}}2\right)+b_n\left(\frac{e^{inx}-e^{-inx}}{2i}\right) =\underbrace{a_0}_{c_0}+\sum_{n=1}^{\infty}\underbrace{\left(\frac{a_n-ib_n}2\right)}_{c_n}e^{inx}+\underbrace{\left(\frac{a_n+ib_n}2\right)}_{c_{-n}}e^{-inx}}
\end{uncoverenv}
\end{frame}

\def\EMtitlesize{26}\def\EMtitleskip{31}
\begin{frame}{صورت مختلط سری فوریه (ادامه)}
\begin{uncoverenv}<2->
\EMeq{c_n =\frac{a_n-ib_n}2 =\frac12\left\{\frac1\pi\int_{2\pi}f(x)\cos nx\,dx-i\frac1\pi\int_{2\pi}f(x)\sin nx\,dx\right\} =\frac1{2\pi}\int_{2\pi}f(x)\big(\cos nx-i\sin nx\big)\,dx \\[4pt] =\frac1{2\pi}\int_{2\pi}f(x)e^{-inx}\,dx}
\end{uncoverenv}
\EMgap
\begin{uncoverenv}<3->
\EMeq{c_{-n} =\frac{a_n+ib_n}2 =\frac1{2\pi}\int_{2\pi}f(x)\big(\cos nx+i\sin nx\big)\,dx =\frac1{2\pi}\int_{2\pi}f(x)e^{inx}\,dx}
\end{uncoverenv}
\EMgap
\begin{EMbox}{imp}{3}{}
\begin{uncoverenv}<4->
\EMeq{f(x) =\sum_{n=-\infty}^{\infty}c_ne^{inx},\qquad c_n =\frac1{2\pi}\int_{2\pi}f(x)e^{-inx}\,dx}
\end{uncoverenv}
\EMgap
\begin{uncoverenv}<5->
\EMp{برای دورۀ تناوب دلخواه:}
\vspace{3bp}
\EMeq{f(x) =\sum_{n=-\infty}^{\infty}c_ne^{i\frac{2\pi}Tnx},\qquad c_n =\frac1T\int_Tf(x)e^{-i\frac{2\pi}Tnx}\,dx}
\end{uncoverenv}
\end{EMbox}
\end{frame}

\def\EMtitlesize{26}\def\EMtitleskip{31}
\begin{frame}{صورت مختلط سری فوریه (ادامه)}
\begin{EMbox}{imp}{1}{}
\begin{uncoverenv}<2->
\EMeq{\frac1T\int_T|f(x)|^2\,dx =\sum_{n=-\infty}^{\infty}|c_n|^2}
\end{uncoverenv}
\end{EMbox}
\end{frame}
